%% ma: L11-B17-0020
%% lop: 11
%% chuong: 5
%% bai: 17
%% dang: 11.17.3
%% loai: TL
%% muc: VD
%% dap_an: "-"
%% nguon: "tai-lieu-hoc-tap-mon-toan-lop-11-theo-chuong-trinh-giao-duc-pho-thong-2018.pdf (Lớp Toán Thầy Duy, Phạm Lê Duy, THPT Chu Văn An), Chương 3 Bài 3, Dạng 3, Ví dụ 20 (tr.393-398)"
%% kien_thuc: "Lớp 11 Bài 17 (hàm số liên tục)"
%% trang_thai: chinh-thuc
%% ghi_chu: "Đã duyệt 03/11/2026."
\begin{ex}
Chứng minh rằng phương trình $a\cos3x+b\cos2x+c\cos x+\sin x=0$ luôn có nghiệm trên $[0;2\pi]$.
\loigiai{\sloppy Đặt $f(x)=a\cos3x+b\cos2x+c\cos x+\sin x$, liên tục trên $[0;2\pi]$. $f(0)=a+b+c$, $f(\pi)=-a+b-c$, $f\left(\dfrac\pi2\right)=-b+1$, $f\left(\dfrac{3\pi}2\right)=-b-1$; tổng bốn giá trị bằng $0$. Nếu có giá trị bằng $0$ thì đó là nghiệm; nếu không thì có hai giá trị trái dấu tại hai điểm thuộc $[0;2\pi]$, nên theo định lí giá trị trung gian phương trình có nghiệm trong $[0;2\pi]$.}
\end{ex}
